\section*{Bloque D. Ampliación: Teoría y Aplicaciones Prácticas (70\% / 30\%)}

\noindent\textit{Preguntas 91 a 180. Ampliación del banco oficial con 63 preguntas teóricas y conceptuales (91 a 153) y 27 aplicaciones cuantitativas prácticas (154 a 180).}

\subsection*{Parte I. Preguntas teóricas y conceptuales (91 a 153)}

\q{¿Por qué el análisis de equilibrio parcial puede resultar insuficiente para evaluar una reforma tributaria de gran escala?}
\begin{opts}
    \item Porque supone que los consumidores nunca maximizan utilidad cuando existen impuestos específicos.
    \item Porque ignora los efectos de retroalimentación sobre precios de bienes relacionados y mercados de factores.
    \item Porque asume que todas las empresas operan bajo rendimientos crecientes a escala en el corto plazo.
    \item Porque impide calcular el excedente del consumidor y la pérdida irrecuperable de eficiencia en mercados lineales.
\end{opts}

\q{En un modelo de equilibrio general con $k$ bienes sin dinero fiduciario, normalizar el precio de un bien como $p_k = 1$ (numerario) es matemáticamente válido porque:}
\begin{opts}
    \item El bien numerario carece de valor de uso y nunca entra en las funciones de utilidad de los agentes.
    \item La Ley de Walras exige que el exceso de demanda del bien numerario sea estrictamente positivo.
    \item Las funciones de exceso de demanda son homogéneas de grado cero en el vector de precios $p$.
    \item Las restricciones presupuestarias de los agentes solo se satisfacen cuando los precios suman uno.
\end{opts}

\q{¿Qué supuesto sobre las preferencias de los consumidores es suficiente para que la Ley de Walras ($\sum_{i=1}^k p_i z_i(p) \equiv 0$) se cumpla como identidad para todo vector de precios $p > 0$?}
\begin{opts}
    \item La no saciedad local, que garantiza que cada agente agote íntegramente su restricción presupuestaria.
    \item La estricta convexidad de las preferencias junto con rendimientos constantes a escala en la producción.
    \item La existencia de preferencias cuasilineales idénticas en todos los consumidores de la economía.
    \item La ausencia de bienes inferiores y el cumplimiento del axioma fuerte de la preferencia revelada.
\end{opts}

\q{En una economía de intercambio puro con $k$ mercados donde se cumple la Ley de Walras, si para un vector de precios estrictamente positivo $p > 0$ hay exactamente $k-1$ mercados en equilibrio ($z_i = 0$), entonces:}
\begin{opts}
    \item El $k$-ésimo mercado puede presentar exceso de oferta o de demanda según la distribución inicial.
    \item El $k$-ésimo mercado presenta necesariamente un exceso de demanda igual a la suma de dotaciones.
    \item Es necesario introducir un bien adicional para cerrar el sistema de ecuaciones del modelo.
    \item El $k$-ésimo mercado también está en equilibrio ($z_k = 0$), dejando $k-1$ ecuaciones.
\end{opts}

\q{Respecto del ámbito de validez de la Ley de Walras ($\sum_{i=1}^k p_i z_i(p) = 0$), es correcto afirmar que:}
\begin{opts}
    \item Solo se cumple exactamente en el vector de precios de equilibrio walrasiano $p^*$ donde $z(p^*) = 0$.
    \item Se cumple como identidad para cualquier vector $p > 0$, esté o no la economía en equilibrio.
    \item Deja de cumplirse si existe un exceso de demanda positivo en el mercado de bienes de consumo.
    \item Únicamente rige en economías con dos bienes y dos consumidores cuyas preferencias son Cobb-Douglas.
\end{opts}

\q{En un equilibrio walrasiano con libre disponibilidad (free disposal), si un bien $j$ presenta un exceso de oferta estricto ($z_j(p^*) < 0$) en equilibrio, entonces su precio de equilibrio $p_j^*$ debe ser:}
\begin{opts}
    \item Estrictamente positivo e igual al costo medio mínimo de producción de la industria.
    \item Negativo, para incentivar a los consumidores a absorber el excedente físico disponible.
    \item Igual a cero ($p_j^* = 0$), actuando como bien libre (holgura complementaria).
    \item Igual al precio del bien numerario para preservar la homogeneidad de las funciones de demanda.
\end{opts}

\q{El proceso de ajuste de precios conocido como tâtonnement walrasiano postula que, fuera del equilibrio:}
\begin{opts}
    \item Un subastador sube precios con exceso de demanda y baja precios con exceso de oferta sin transar antes del equilibrio.
    \item Los agentes realizan intercambios bilaterales inmediatos a precios de desequilibrio modificando sus dotaciones en cada ronda de negociación.
    \item Las empresas ajustan únicamente sus cantidades producidas mientras el nivel relativo de precios permanece fijo durante todo el ajuste.
    \item El gobierno fija cuotas obligatorias de consumo hasta que las tasas marginales de sustitución se igualen entre todos los agentes.
\end{opts}

\q{Si en una economía real de equilibrio general todos los precios monetarios se multiplican por $5$ (sin saldos monetarios reales ni ilusión monetaria), el efecto sobre las asignaciones de equilibrio es:}
\begin{opts}
    \item Un aumento proporcional del consumo de los bienes normales y una caída de los bienes inferiores.
    \item Un desplazamiento de la curva de contrato hacia el origen del agente con menor dotación inicial.
    \item Una redistribución del ingreso real desde los oferentes netos hacia los demandantes netos de cada bien.
    \item Nulo: las cantidades demandadas, ofrecidas y las utilidades de equilibrio permanecen inalteradas.
\end{opts}

\q{En una Caja de Edgeworth de intercambio puro entre dos consumidores $A$ y $B$ con dotaciones iniciales $\omega_A = (\omega_{1A}, \omega_{2A})$ y $\omega_B = (\omega_{1B}, \omega_{2B})$, el ancho y la altura de la caja representan:}
\begin{opts}
    \item Las cantidades iniciales que posee el consumidor $A$ y las que posee el consumidor $B$, respectivamente.
    \item Las dotaciones agregadas totales: $X = \omega_{1A} + \omega_{1B}$ y $Y = \omega_{2A} + \omega_{2B}$.
    \item Los niveles máximos de utilidad que pueden alcanzar los agentes $A$ y $B$ sobre la frontera.
    \item Los precios relativos de equilibrio del bien $1$ y del bien $2$ expresados en términos del numerario.
\end{opts}

\q{Si en la dotación inicial $W$ de una Caja de Edgeworth se cumple $TMS_{1,2}^A(W) > TMS_{1,2}^B(W)$ (con preferencias estrictamente convexas), entonces existe un intercambio mutuamente beneficioso en el cual:}
\begin{opts}
    \item El consumidor $A$ entrega unidades del bien $1$ al consumidor $B$ a cambio de unidades del bien $2$.
    \item Ambos consumidores reducen su consumo del bien $1$ para incrementar la dotación total de la economía.
    \item El consumidor $A$ recibe unidades del bien $1$ de $B$ y le entrega unidades del bien $2$.
    \item Ningún consumidor desea intercambiar porque $W$ ya pertenece automáticamente al núcleo.
\end{opts}

\q{En una Caja de Edgeworth donde ambos consumidores tienen preferencias estrictamente crecientes pero diferenciables solo en el interior, una asignación en el borde de la caja (solución de esquina) puede ser Pareto eficiente:}
\begin{opts}
    \item Incluso si $TMS^A \neq TMS^B$, cuando ningún movimiento factible en la caja mejora a ambos.
    \item Únicamente cuando se cumple la igualdad estricta $TMS^A = TMS^B = 1$ en uno de los vértices de la caja.
    \item Nunca, porque toda asignación eficiente en el sentido de Pareto debe ser estrictamente interior.
    \item Solo si ambos consumidores poseen exactamente la misma función de utilidad Cobb-Douglas simétrica.
\end{opts}

\q{¿Cuál es la diferencia precisa entre la Curva de Contrato y el Núcleo (Core) en una economía de intercambio puro $2\times 2$ con dotación inicial $W$?}
\begin{opts}
    \item La Curva de Contrato depende de la dotación inicial $W$, mientras el Núcleo es independiente de $W$.
    \item El Núcleo incluye todas las asignaciones del lente de mejora, aunque no sean eficientes en sentido de Pareto.
    \item Ambos conjuntos coinciden siempre exactamente para cualquier dotación inicial ubicada en el interior de la caja.
    \item El Núcleo es el subconjunto de la Curva de Contrato donde ambos agentes alcanzan al menos su utilidad en $W$.
\end{opts}

\q{En una economía de intercambio puro con preferencias localmente no saciadas, toda asignación de equilibrio walrasiano pertenece necesariamente al Núcleo porque:}
\begin{opts}
    \item El subastador prohíbe legalmente que los consumidores formen coaliciones para bloquear asignaciones.
    \item Ninguna coalición puede mejorar a sus miembros con sus dotaciones sin violar presupuestos.
    \item Todas las asignaciones del núcleo presentan exactamente el mismo vector de precios relativos de equilibrio.
    \item Los consumidores obtienen exactamente la misma utilidad absoluta en el equilibrio walrasiano.
\end{opts}

\q{Según el teorema de límite del núcleo (Edgeworth / Debreu-Scarf), al replicar una economía de intercambio aumentando indefinidamente el número de consumidores de cada tipo ($r \to \infty$):}
\begin{opts}
    \item El Núcleo se expande hasta cubrir toda la Caja de Edgeworth y desaparece el equilibrio walrasiano.
    \item El Núcleo se vuelve vacío debido a la inestabilidad de las coaliciones de bloqueo.
    \item El Núcleo se contrae progresivamente hasta coincidir con el conjunto de equilibrios walrasianos.
    \item La Curva de Contrato se transforma en una línea recta independiente de las preferencias individuales.
\end{opts}

\q{Si en una Caja de Edgeworth ambos consumidores tienen preferencias de complementos perfectos $U_A = \min(x_A, y_A)$ y $U_B = \min(x_B, y_B)$, y la dotación total cumple $X > Y$, entonces el conjunto de asignaciones Pareto eficientes:}
\begin{opts}
    \item Es una banda bidimensional donde $Y$ se asigna en proporción $1:1$ y sobra el bien $X$.
    \item Se reduce a un único punto en el centro geométrico de la caja independientemente de $X$ e $Y$.
    \item Es exactamente la diagonal principal que une $O_A$ con $O_B$ sin ningún espesor lateral.
    \item Es vacío porque las funciones de utilidad tipo Leontief no son diferenciables en los vértices.
\end{opts}

\q{Si ambos consumidores en una Caja de Edgeworth tienen preferencias cuasilineales en el bien $2$ ($U_i = v_i(x_{1i}) + x_{2i}$, con $v_i'' < 0$), entonces para las asignaciones interiores eficientes:}
\begin{opts}
    \item La asignación del bien $2$ queda unívocamente fija y solo varía la distribución del bien $1$.
    \item La Curva de Contrato es siempre una línea recta con pendiente igual a $45^\circ$ desde $O_A$.
    \item El precio relativo de equilibrio depende exclusivamente de cómo se distribuya el bien $2$ inicialmente.
    \item La asignación del bien $1$ es única e independiente del bien $2$ (Curva de Contrato vertical).
\end{opts}

\q{En la Caja de Edgeworth, la curva de oferta (o curva precio-consumo) de un consumidor trazada desde la dotación inicial $W$ representa:}
\begin{opts}
    \item El conjunto de todas las asignaciones eficientes en el sentido de Pareto dentro de la caja.
    \item Las canastas óptimas del agente conforme varía el precio relativo $p_1 / p_2$ por $W$.
    \item La frontera máxima de producción que puede generar el consumidor combinando ambos bienes.
    \item El conjunto de asignaciones que dejan al consumidor exactamente indiferente frente a la dotación $W$.
\end{opts}

\q{En una economía con dos factores ($L, K$) y dos bienes ($X, Y$), la eficiencia en la asignación de factores (estar sobre la FPP) exige en el interior de la Caja de Edgeworth de producción que:}
\begin{opts}
    \item El producto marginal del trabajo sea idéntico en ambas industrias ($PMg_L^X = PMg_L^Y$).
    \item La relación capital-trabajo sea exactamente igual en la producción de ambos bienes ($K_X/L_X = K_Y/L_Y$).
    \item Las tasas marginales de sustitución técnica se igualen entre industrias ($TMST_{L,K}^X = TMST_{L,K}^Y$).
    \item El costo medio de producir el bien $X$ iguale al costo medio de producir el bien $Y$.
\end{opts}

\q{La Tasa Marginal de Transformación ($TMT_{X,Y} = -\frac{dY}{dX}$) a lo largo de la Frontera de Posibilidades de Producción es económicamente igual a:}
\begin{opts}
    \item El cociente de costos marginales $\frac{CMg_X}{CMg_Y}$ (costo de oportunidad de $X$).
    \item El cociente de costos medios totales $\frac{CMe_X}{CMe_Y}$ evaluado en la escala mínima eficiente.
    \item La razón entre la dotación total de capital y la dotación total de trabajo de la economía.
    \item El cociente inverso de los productos marginales del trabajo multiplicado por el salario real.
\end{opts}

\q{Incluso cuando ambos sectores operan con rendimientos constantes a escala, la Frontera de Posibilidades de Producción (FPP) es estrictamente cóncava hacia el origen si:}
\begin{opts}
    \item Las dos industrias utilizan capital y trabajo exactamente en la misma proporción factorial.
    \item Solo existe un único factor productivo homogéneo que puede trasladarse sin costo entre sectores.
    \item Los consumidores presentan preferencias estrictamente convexas y homogéneas de grado uno.
    \item Las dos industrias tienen intensidades factoriales distintas ($X$ intensivo en $L$, $Y$ en $K$).
\end{opts}

\q{¿Bajo qué conjunto de condiciones la Frontera de Posibilidades de Producción (FPP) de una economía $2\times 2$ resulta ser una línea recta?}
\begin{opts}
    \item Cuando ambos bienes presentan rendimientos crecientes a escala con intensidades factoriales distintas.
    \item Cuando hay RCE y ambos sectores usan factores en idéntica proporción (o un solo factor).
    \item Cuando las preferencias de todos los consumidores son de sustitutos perfectos con $TMS$ constante.
    \item Cuando existe desempleo involuntario de trabajo y subutilización de capital en ambos sectores.
\end{opts}

\q{Si una economía se encuentra sobre su FPP (eficiencia técnica) y sobre la Curva de Contrato de consumo ($TMS^A = TMS^B$), pero $TMS_{X,Y} > TMT_{X,Y}$, entonces:}
\begin{opts}
    \item La asignación ya es un Óptimo de Pareto global porque no existe desempleo de factores.
    \item Se está produciendo demasiado del bien $X$ y muy poco del bien $Y$ respecto del óptimo social.
    \item Se puede mejorar el bienestar produciendo más del bien $X$ y menos del bien $Y$.
    \item El precio relativo de mercado del bien $X$ es necesariamente mayor que su valoración marginal.
\end{opts}

\q{En un equilibrio general competitivo con producción, el sistema de precios descentralizado logra simultáneamente las tres condiciones de eficiencia paretiana porque:}
\begin{opts}
    \item Cada agente toma precios, igualando $TMS = p_X/p_Y = CMg_X/CMg_Y = TMT$ y $TMST = w/r$.
    \item El planificador central fija salarios y tasas de interés distintos para cada sector según su productividad.
    \item Todas las firmas igualan su costo medio al precio y todos los consumidores reciben idéntico ingreso.
    \item La Ley de Walras obliga a que el producto marginal del capital iguale a la tasa marginal de sustitución.
\end{opts}

\q{En la economía de Robinson Crusoe con un bien de consumo $c$ y ocio $l$ (donde el trabajo es $L = T - l$), la condición de eficiencia paretiana interior es:}
\begin{opts}
    \item El producto medio del trabajo debe igualar a la utilidad marginal del consumo.
    \item El consumo total debe ser igual al nivel de ocio independientemente de la tecnología.
    \item La tasa marginal de sustitución entre ocio y consumo debe ser igual al costo fijo total.
    \item La tasa marginal de sustitución ocio-consumo iguala al $PMg_L$ ($TMS_{l,c} = PMg_L$).
\end{opts}

\q{Si una industria opera con rendimientos crecientes a escala (conjunto de producción no convexo), el equilibrio competitivo descentralizado enfrenta un problema fundamental porque:}
\begin{opts}
    \item Las empresas obtienen beneficios económicos infinitos cuando igualan el precio al costo marginal.
    \item Fijar $P = CMg$ genera pérdidas ($CMg < CMe$), impidiendo un equilibrio competitivo.
    \item La Frontera de Posibilidades de Producción se desplaza hacia el origen eliminando todo consumo.
    \item La Ley de Walras deja de cumplirse cuando los costos marginales son decrecientes.
\end{opts}

\q{Si en una economía competitiva con producción se introduce un impuesto ad-valorem únicamente sobre el consumo del bien $X$ (devolviendo la recaudación como suma alzada), la eficiencia global se rompe porque:}
\begin{opts}
    \item La economía pasa a operar estrictamente en el interior de su Frontera de Posibilidades de Producción.
    \item Se viola la eficiencia técnica en el uso de factores entre sectores ($TMST^X \neq TMST^Y$).
    \item Se genera una brecha entre $TMS = p_X^c / p_Y$ y $TMT = p_X^p / p_Y$, rompiendo $TMS = TMT$.
    \item Los consumidores dejan de igualar sus tasas marginales de sustitución entre sí ($TMS^A \neq TMS^B$).
\end{opts}

\q{Respecto de los supuestos necesarios para que se cumpla el Primer Teorema Fundamental del Bienestar, es correcto señalar que:}
\begin{opts}
    \item Requiere no saciedad local, mercados completos y toma de precios, pero no requiere convexidad.
    \item Exige estricta convexidad de preferencias y conjuntos de producción para garantizar la eficiencia.
    \item Solo se cumple si la distribución inicial de dotaciones es equitativa entre todos los agentes.
    \item Exige que todas las empresas tengan rendimientos decrecientes a escala y beneficios nulos.
\end{opts}

\q{Una limitación normativa central del Primer Teorema del Bienestar es que:}
\begin{opts}
    \item Demuestra que el equilibrio competitivo maximiza el bienestar de la función social rawlsiana.
    \item Garantiza que las asignaciones competitivas eliminan las desigualdades extremas de dotación.
    \item Implica que toda intervención del Estado reduce necesariamente el bienestar social agregado.
    \item Una asignación puede ser Pareto eficiente y, a la vez, sumamente desigual e inequitativa.
\end{opts}

\q{¿Por qué el Segundo Teorema Fundamental del Bienestar requiere que las preferencias y los conjuntos de producción sean convexos?}
\begin{opts}
    \item Para asegurar que la Ley de Walras se cumpla en todos los mercados simultáneamente.
    \item Para aplicar el hiperplano separador y hallar precios que sostengan el óptimo deseado.
    \item Para garantizar que el equilibrio competitivo resultante sea único y globalmente estable.
    \item Para evitar que los consumidores prefieran el ocio frente al consumo de bienes materiales.
\end{opts}

\q{Según el Segundo Teorema del Bienestar, para alcanzar cualquier asignación Pareto eficiente deseada como equilibrio competitivo, el Estado debe emplear:}
\begin{opts}
    \item Controles de precios máximos en los bienes básicos y pisos salariales en el mercado laboral.
    \item Impuestos específicos sobre los bienes de lujo combinados con subsidios a los bienes de primera necesidad.
    \item Transferencias de suma alzada (lump-sum) que redistribuyan riqueza sin alterar los precios relativos.
    \item Cuotas físicas obligatorias de producción asignadas centralmente a cada empresa de la economía.
\end{opts}

\q{En la práctica de la política económica, la dificultad principal para implementar transferencias de suma alzada (lump-sum) como sugiere el Segundo Teorema radica en que:}
\begin{opts}
    \item El Estado no observa habilidades innatas, por lo que los impuestos terminan distorsionando el esfuerzo.
    \item Las transferencias de suma alzada siempre rompen la igualdad de las tasas marginales de sustitución entre consumidores.
    \item La teoría económica demuestra que toda transferencia monetaria genera una pérdida irrecuperable mayor que un impuesto al consumo.
    \item El teorema del hiperplano separador únicamente es válido cuando hay exactamente dos bienes en la economía.
\end{opts}

\q{¿Cuál de los siguientes escenarios invalida directamente el Primer Teorema del Bienestar al impedir que el equilibrio de mercado sea un Óptimo de Pareto?}
\begin{opts}
    \item Una distribución inicial altamente concentrada de las dotaciones de capital en pocos agentes.
    \item La existencia de preferencias heterogéneas donde cada consumidor valora de forma distinta los bienes.
    \item La presencia de precios relativos diferentes de la unidad en el equilibrio competitivo.
    \item La existencia de externalidades no internalizadas, poder de mercado o mercados incompletos.
\end{opts}

\q{De acuerdo con el Teorema del Segundo Mejor (Lipsey y Lancaster), si en una economía existe una distorsión inamovible en un mercado (por ejemplo, un monopolio no regulable), entonces:}
\begin{opts}
    \item Eliminar las distorsiones en todos los demás mercados garantiza siempre una mejora del bienestar social.
    \item Cumplir $P = CMg$ en el resto de mercados puede no ser socialmente óptimo (second-best).
    \item La economía alcanza automáticamente el Óptimo de Pareto mediante la Ley de Walras.
    \item El Segundo Teorema del Bienestar permite restaurar el óptimo global sin necesidad de información.
\end{opts}

\q{Se afirma que el Segundo Teorema del Bienestar permite separar teóricamente el problema de la eficiencia del problema de la equidad porque:}
\begin{opts}
    \item El mercado competitivo produce automáticamente la distribución más equitativa de utilidades.
    \item Las transferencias distorsionadoras pueden igualar utilidades sin reducir nunca la producción agregada.
    \item Cualquier meta distributiva eficiente se logra redistribuyendo dotaciones y dejando operar al mercado.
    \item La Frontera de Posibilidades de Utilidad es siempre una recta horizontal independiente de los precios.
\end{opts}

\q{En el marco de equilibrio general con incertidumbre (Arrow-Debreu), para que el Primer Teorema del Bienestar siga siendo válido se requiere que existan mercados completos, es decir:}
\begin{opts}
    \item Que exista un mercado competitivo para cada bien en cada estado de la naturaleza (contingentes).
    \item Que todos los agentes sean neutrales al riesgo y posean exactamente la misma probabilidad subjetiva.
    \item Que el gobierno elimine toda fuente de incertidumbre climática o tecnológica mediante regulación directa.
    \item Que los precios de los bienes sean idénticos en todos los estados futuros de la naturaleza.
\end{opts}

\q{Para demostrar matemáticamente la existencia de un equilibrio walrasiano en el modelo de Arrow-Debreu cuando las demandas son funciones continuas (o correspondencias), se utiliza:}
\begin{opts}
    \item El Teorema de la Envolvente combinado con la identidad de Roy en equilibrio parcial.
    \item El Teorema de Coase sobre derechos de propiedad y negociación sin costos de transacción.
    \item La desigualdad de Jensen aplicada a la función de beneficios agregada de las firmas.
    \item El Teorema de Punto Fijo de Brouwer (funciones) o de Kakutani (correspondencias).
\end{opts}

\q{El resultado fundamental del Teorema de Sonnenschein-Mantel-Debreu (SMD) sobre la función de exceso de demanda agregada $z(p)$ establece que:}
\begin{opts}
    \item Las preferencias individuales bien comportadas garantizan siempre un único equilibrio walrasiano globalmente estable.
    \item Casi cualquier función continua, homogénea de grado cero y que cumpla la Ley de Walras puede ser un $z(p)$ agregado.
    \item El equilibrio general no existe cuando hay más de tres bienes y dos consumidores en la economía.
    \item La función de exceso de demanda agregada hereda siempre la propiedad de monotonicidad del axioma débil.
\end{opts}

\q{¿Qué condición adicional sobre la función de exceso de demanda agregada $z(p)$ garantiza la unicidad y estabilidad global del equilibrio walrasiano?}
\begin{opts}
    \item Que todos los bienes sean complementarios brutos entre sí ($\partial z_i / \partial p_j < 0$ para todo $j \neq i$).
    \item Que las preferencias de todos los consumidores sean lexicográficas y no diferenciables.
    \item Que todos los bienes sean sustitutos brutos ($\partial z_i / \partial p_j > 0$ para todo $j \neq i$).
    \item Que la economía tenga un número par de bienes y un número impar de consumidores.
\end{opts}

\q{Si la función de exceso de demanda agregada $z(p)$ satisface el Axioma Débil de la Preferencia Revelada (WARP), entonces el conjunto de equilibrios walrasianos:}
\begin{opts}
    \item Es convexo y, en economías regulares, el equilibrio relativo de precios es único.
    \item Es necesariamente vacío debido a la contradicción con la Ley de Walras.
    \item Contiene siempre un número impar mayor o igual a tres equilibrios aislados.
    \item Depende del orden en que el subastador anuncia los precios durante el tâtonnement.
\end{opts}

\q{En una economía de intercambio puro regular (con matriz jacobiana de exceso de demanda no singular en equilibrio), el Teorema del Índice (Dierker) implica que el número de equilibrios walrasianos es:}
\begin{opts}
    \item Siempre igual a uno, independientemente de los efectos ingreso en la economía.
    \item Siempre un número par mayor o igual a dos, salvo cuando las preferencias son Cobb-Douglas.
    \item Infinito no numerable en cualquier economía con al menos dos bienes y dos consumidores.
    \item Finito y necesariamente un número impar ($1, 3, 5, \dots$).
\end{opts}

\q{En una economía de intercambio puro con dos bienes (normalizando $p_2 = 1$), el equilibrio walrasiano $p_1^*$ es localmente estable bajo el tâtonnement $\dot{p}_1 = \alpha z_1(p_1, 1)$ si en $p_1^*$ se cumple:}
\begin{opts}
    \item $\frac{dz_1}{dp_1}(p_1^*) > 0$, es decir, el exceso de demanda crece al subir el precio.
    \item $\frac{dz_1}{dp_1}(p_1^*) < 0$, pues subir $p_1$ reduce el exceso de demanda del bien $1$.
    \item $\frac{dz_1}{dp_1}(p_1^*) = 0$, garantizando que el precio permanezca constante ante perturbaciones.
    \item $z_1(p_1^*, 1) > 0$, manteniendo una presión permanente al alza sobre el precio relativo.
\end{opts}

\q{¿Qué fenómeno microeconómico puede provocar que $z_1(p_1, 1)$ tenga pendiente positiva en torno a un equilibrio walrasiano, generando equilibrios múltiples o inestables?}
\begin{opts}
    \item El efecto sustitución puro, que siempre empuja la demanda en la misma dirección que el precio.
    \item La existencia de costos de transacción nulos en la negociación entre consumidores.
    \item Un efecto ingreso o riqueza sobre la dotación que domine al efecto sustitución.
    \item La homogeneidad de grado cero de las funciones de demanda marshallianas individuales.
\end{opts}

\q{En el modelo intertemporal de Arrow-Debreu, un mismo bien físico (por ejemplo, agua potable) entregado en dos fechas distintas o en dos ubicaciones geográficas distintas se modela como:}
\begin{opts}
    \item Dos bienes distintos con precios propios (análisis intertemporal y espacial).
    \item Un único bien homogéneo con un solo precio relativo idéntico en todos los períodos.
    \item Una externalidad negativa intertemporal que invalida el Primer Teorema del Bienestar.
    \item Un bien público puro excluido de las restricciones presupuestarias individuales.
\end{opts}

\q{En la demostración de existencia del equilibrio walrasiano, el comportamiento de frontera (boundary behavior) establece que cuando el precio de un bien deseable tiende a cero ($p_j \to 0$):}
\begin{opts}
    \item El exceso de demanda de dicho bien tiende a menos infinito por saturación de los agentes.
    \item La utilidad de todos los consumidores converge a cero debido a la caída del ingreso nominal.
    \item El equilibrio se desplaza necesariamente hacia una solución de esquina con $p_j^* = 0$.
    \item La norma del exceso de demanda tiende a infinito, empujando $p^*$ al interior ($p^* > 0$).
\end{opts}

\q{En un modelo de equilibrio general $2\times 2$ con rendimientos constantes a escala y producción diversificada donde $X$ es intensivo en trabajo ($L$) e $Y$ es intensivo en capital ($K$), el Teorema de Stolper-Samuelson establece que un aumento en el precio relativo $p_X / p_Y$:}
\begin{opts}
    \item Eleva tanto el salario real como la renta real del capital en la misma proporción que $p_X$.
    \item Eleva el salario real y reduce la renta real del capital en términos de ambos bienes.
    \item Reduce el salario real debido a la mayor inflación provocada por el encarecimiento del bien $X$.
    \item Deja inalteradas las remuneraciones reales de ambos factores al mantenerse constantes las dotaciones.
\end{opts}

\q{En el modelo $2\times 2$ con precios de bienes constantes y pleno empleo de factores, si $X$ es intensivo en trabajo ($L$) e $Y$ es intensivo en capital ($K$), un aumento exógeno en la dotación de trabajo ($\bar{L}$) provoca:}
\begin{opts}
    \item Un aumento proporcional idéntico en la producción de ambos bienes $X$ e $Y$.
    \item Una caída en la producción de $X$ y un aumento en la producción de $Y$ para absorber el capital.
    \item Una expansión más que proporcional de $X$ y una contracción absoluta de la producción de $Y$.
    \item Una reducción del salario real acompañada de un aumento en el precio relativo del bien $X$.
\end{opts}

\q{El Teorema de Igualación de Precios de Factores (Samuelson) en el modelo Heckscher-Ohlin $2\times 2$ implica que el libre comercio de bienes entre dos países con igual tecnología y producción diversificada:}
\begin{opts}
    \item Iguala los salarios y las rentas del capital entre ambos países aun sin movilidad internacional de factores.
    \item Exige necesariamente la migración libre de trabajadores y capitales para igualar los precios de los bienes.
    \item Hace que el país abundante en trabajo pague salarios permanentemente mayores que el país abundante en capital.
    \item Elimina toda ventaja comparativa y lleva el volumen de comercio internacional a cero en el largo plazo.
\end{opts}

\q{Según el Teorema de Heckscher-Ohlin en equilibrio general internacional $2\times 2$, un país relativamente abundante en capital tenderá a exportar:}
\begin{opts}
    \item El bien que tenga el menor costo marginal en términos de trabajo físico por unidad.
    \item Ambos bienes simultáneamente para acumular reservas internacionales en el largo plazo.
    \item El bien cuya demanda interna sea más inelástica respecto del ingreso nacional.
    \item El bien que utiliza el capital de manera relativamente intensiva en su producción.
\end{opts}

\q{En el modelo de equilibrio general de Harberger para analizar un impuesto al capital aplicado únicamente en el sector corporativo ($X$), la incidencia del tributo se descompone en:}
\begin{opts}
    \item Un efecto renta ricardiano y un efecto de externalidad pecuniaria sobre la curva de oferta.
    \item Un efecto sustitución de factores y un efecto producto según la intensidad factorial de $X$.
    \item Un efecto de demanda agregada y un efecto de ilusión monetaria en los salarios nominales.
    \item Un efecto de competencia imperfecta y un efecto de racionamiento por colas en el sector $Y$.
\end{opts}

\q{En el modelo de Harberger, si el sector gravado con un impuesto al capital ($X$) es muy intensivo en trabajo y la elasticidad de sustitución entre factores es muy baja, entonces el impuesto al capital en $X$ puede:}
\begin{opts}
    \item Ser absorbido íntegramente por los propietarios del capital en ambos sectores sin afectar al salario.
    \item Aumentar la producción del sector gravado $X$ por encima de su nivel inicial sin impuestos.
    \item Terminar elevando la remuneración relativa del capital ($r/w$) y recayendo sobre los trabajadores.
    \item Eliminar por completo la pérdida irrecuperable de eficiencia en toda la economía.
\end{opts}

\q{En equilibrio general con producción, la relación de precios $p_X/p_Y$ entre dos sectores queda determinada simultáneamente por:}
\begin{opts}
    \item La tangencia entre la FPP y las preferencias sociales (oferta relativa y demanda relativa).
    \item Únicamente los costos medios mínimos de producción, sin ninguna influencia de la demanda.
    \item Únicamente las preferencias de los consumidores, sin influencia de las dotaciones factoriales.
    \item La razón entre el número de empresas en el sector $X$ y el número de empresas en el sector $Y$.
\end{opts}

\q{La Frontera de Posibilidades de Utilidad (FPU) representa geométricamente:}
\begin{opts}
    \item Las combinaciones máximas de bienes $X$ e $Y$ que la economía puede producir con sus factores.
    \item Todas las combinaciones de utilidades alcanzables mediante asignaciones dentro del lente de mejora.
    \item El conjunto de precios relativos que vacían simultáneamente todos los mercados competitivos.
    \item La imagen en el espacio de utilidades $(U_A, U_B)$ de todos los Óptimos de Pareto de la economía.
\end{opts}

\q{Una crítica formal al criterio de compensación de Kaldor-Hicks es la paradoja de reversión de Scitovsky, que ocurre cuando:}
\begin{opts}
    \item Una política genera beneficios positivos para todos los individuos sin necesidad de compensación.
    \item El paso de $A$ a $B$ cumple Kaldor-Hicks y, a la vez, el regreso de $B$ a $A$ también lo cumple.
    \item Los ganadores compensan efectivamente a los perdedores logrando una mejora estricta de Pareto.
    \item La Frontera de Posibilidades de Utilidad tiene pendiente negativa en todo su recorrido.
\end{opts}

\q{En una familia de funciones de bienestar social CES $W = (U_1^\rho + U_2^\rho)^{1/\rho}$ con $\rho \le 1$, el grado de aversión social a la desigualdad aumenta conforme:}
\begin{opts}
    \item Aumenta $\rho$ hacia $+\infty$, acercándose al criterio utilitarista puro de Bentham.
    \item Se mantiene $\rho = 1$, pues las curvas de isobienestar en forma de recta penalizan la desigualdad.
    \item Disminuye $\rho$ desde $1$ (Bentham) pasando por $0$ (Nash) hacia $-\infty$ (Rawls / maximin).
    \item Se igualan las utilidades marginales del ingreso a uno mediante impuestos ad-valorem.
\end{opts}

\q{El Teorema de Imposibilidad de Arrow demuestra que, con al menos tres alternativas sociales, no existe ninguna función de bienestar social (regla de agregación de rankings ordinales) que satisfaga simultáneamente:}
\begin{opts}
    \item Dominio universal, eficiencia de Pareto, independencia de alternativas irrelevantes (IIA) y no dictadura.
    \item Dominio restringido a preferencias unimodales, regla de mayoría simple y ausencia de ciclos de Condorcet.
    \item Comparabilidad interpersonal cardinal de utilidades, eficiencia productiva y equidad rawlsiana.
    \item No saciedad local, convexidad estricta de preferencias, mercados completos y Ley de Walras.
\end{opts}

\q{El axioma de Independencia de Alternativas Irrelevantes (IIA) en el Teorema de Arrow exige que el ordenamiento social entre dos opciones $x$ e $y$:}
\begin{opts}
    \item Dependa de la intensidad cardinal con la que cada votante prefiere $x$ sobre $y$.
    \item Cambie siempre que se introduzca una tercera alternativa $z$ que todos los votantes consideran inferior.
    \item Coincida necesariamente con la preferencia individual del votante con mayor ingreso de la sociedad.
    \item Dependa solo de las preferencias individuales entre $x$ e $y$, sin influirse por terceros $z$.
\end{opts}

\q{¿Qué restricción sobre las preferencias individuales permite superar la Paradoja de Condorcet y garantiza que la regla de mayoría simple produzca un ganador único (el preferido por el votante mediano)?}
\begin{opts}
    \item Que todos los votantes tengan funciones de utilidad idénticas y estrictamente lineales.
    \item Que las preferencias de todos sean unimodales (single-peaked) en una dimensión.
    \item Que el número de alternativas sometidas a votación sea estrictamente igual a dos.
    \item Que los votantes puedan comprar y vender votos en un mercado competitivo sin fricciones.
\end{opts}

\q{Geométricamente, el óptimo de bienestar social (bliss point) bajo una función de bienestar social diferenciable $W(U_A, U_B)$ se ubica en el punto donde:}
\begin{opts}
    \item La recta de $45^\circ$ intersecta a la Frontera de Posibilidades de Producción en el espacio de bienes.
    \item La curva de isobienestar social más baja posible corta a la FPU con pendiente positiva.
    \item La curva de isobienestar más alta alcanzable es tangente a la Frontera de Posibilidades de Utilidad.
    \item La utilidad del individuo $A$ alcanza su valor máximo absoluto haciendo cero la utilidad de $B$.
\end{opts}

\q{En la regulación tarifaria de agua potable (SUNASS), la recomendación teórica derivada del Segundo Teorema del Bienestar para proteger a hogares de bajos ingresos con mínima pérdida de eficiencia es:}
\begin{opts}
    \item Cobrar tarifas basadas en costo marginal y asistir a hogares pobres con bonos de suma alzada.
    \item Fijar una tarifa igual a cero para todos los usuarios residenciales financiando a la EPS con deuda externa.
    \item Aumentar artificialmente el precio del agua a las industrias para subsidiar el consumo ilimitado residencial.
    \item Racionar físicamente el servicio por horas iguales para todos los distritos sin diferenciar ingresos.
\end{opts}

\q{Cuando el Estado evalúa la viabilidad social de un proyecto de planta de tratamiento de aguas residuales mediante Análisis Costo-Beneficio ($VAN\text{ social} > 0$), el criterio normativo subyacente es:}
\begin{opts}
    \item El criterio estricto de Pareto, que exige que ningún contribuyente pierda un solo sol con el proyecto.
    \item El criterio rawlsiano puro, que solo contabiliza los beneficios recibidos por el hogar más pobre.
    \item El Teorema de Imposibilidad de Arrow, que prohíbe agregar beneficios de distintos distritos.
    \item El criterio de Kaldor-Hicks: beneficios totales superan costos (compensación potencial).
\end{opts}

\q{Si un planificador de política hídrica decide priorizar el presupuesto público para llevar agua potable segura primero a las comunidades rurales sin acceso (el grupo peor situado), antes que mejorar la presión en zonas urbanas ricas, su criterio refleja:}
\begin{opts}
    \item Una función de bienestar social utilitarista pura que maximiza únicamente la suma monetaria de excedentes.
    \item Una función de bienestar social rawlsiana (maximin), priorizando al grupo peor situado.
    \item El principio de neutralidad fiscal del Primer Teorema del Bienestar sin intervención estatal.
    \item La paradoja de Condorcet aplicada a la asignación de inversiones en infraestructura sanitaria.
\end{opts}

\q{En una cuenca sujeta a sequías severas donde no existen seguros climáticos ni contratos contingentes de agua entre usuarios urbanos y agrícolas, el Primer Teorema del Bienestar falla principalmente por:}
\begin{opts}
    \item La violación del supuesto de no saciedad local en el consumo de agua potable.
    \item La existencia de rendimientos constantes a escala en la agricultura de riego.
    \item Mercados incompletos, que impiden asignar eficientemente el riesgo hidrológico.
    \item La imposibilidad de definir un bien numerario en presencia de variaciones estacionales.
\end{opts}

\q{En el sector de agua potable y alcantarillado, el equilibrio de mercado no regulado no alcanza el Óptimo de Pareto del Primer Teorema debido fundamentalmente a:}
\begin{opts}
    \item Monopolio natural en redes (no convexidad) junto con externalidades sanitarias y ambientales.
    \item La existencia de demasiadas empresas competidoras precio-aceptantes ofreciendo el mismo servicio.
    \item El cumplimiento estricto de la Ley de Walras en los mercados de factores productivos locales.
    \item La perfecta información que poseen los usuarios sobre la calidad bacteriológica del agua.
\end{opts}

\subsection*{Parte II. Aplicaciones prácticas y ejercicios cuantitativos (154 a 180)}

\q{En una economía de intercambio puro con tres bienes donde se cumple la Ley de Walras, el vector de precios es $p = (2, 3, 5)$ y los excesos de demanda de los dos primeros bienes son $z_1 = 10$ y $z_2 = -5$. ¿Cuál es el exceso de demanda $z_3$?}
\begin{opts}
    \item $z_3 = -1$ (exceso de oferta)
    \item $z_3 = +2$ (exceso de demanda)
    \item $z_3 = -3$ (exceso de oferta)
    \item $z_3 = +5$ (exceso de demanda)
\end{opts}

\q{Una economía de dos bienes presenta la función de exceso de demanda $z_1(p_1, p_2) = \frac{6p_2 - 2p_1}{p_1}$. Si se fija el bien $2$ como numerario ($p_2 = 1$), el precio de equilibrio $p_1^*$ y el grado de homogeneidad de $z_1$ son:}
\begin{opts}
    \item $p_1^* = 3$ y grado $1$
    \item $p_1^* = 3$ y grado $0$
    \item $p_1^* = 2$ y grado $0$
    \item $p_1^* = 6$ y grado $0$
\end{opts}

\q{En una Caja de Edgeworth con dotaciones totales $X = 12$ e $Y = 12$, las utilidades son $U_A = x_A^{1/3} y_A^{2/3}$ y $U_B = x_B^{2/3} y_B^{1/3}$. Si una asignación Pareto eficiente tiene $x_A = 4$, ¿cuál es el valor correspondiente de $y_A$?}
\begin{opts}
    \item $y_A = 4$
    \item $y_A = 6$
    \item $y_A = 8$
    \item $y_A = 9$
\end{opts}

\q{Considere una economía de intercambio con $U_A = x_A y_A$, $U_B = x_B y_B$ y dotaciones iniciales $W_A = (12, 0)$ y $W_B = (0, 6)$. Con $p_y = 1$, el precio relativo de equilibrio $p_x^*$ y la canasta óptima de $A$ son:}
\begin{opts}
    \item $p_x^* = 2.0$ y $(x_A^*, y_A^*) = (6, 6)$
    \item $p_x^* = 1.0$ y $(x_A^*, y_A^*) = (6, 3)$
    \item $p_x^* = 0.5$ y $(x_A^*, y_A^*) = (3, 6)$
    \item $p_x^* = 0.5$ y $(x_A^*, y_A^*) = (6, 3)$
\end{opts}

\q{En una economía $2\times 2$ con $U_A = x_A^{1/4} y_A^{3/4}$, $U_B = x_B^{3/4} y_B^{1/4}$ y dotaciones $W_A = (4, 4)$ y $W_B = (4, 12)$, ¿cuál es el precio relativo de equilibrio $p_x / p_y$?}
\begin{opts}
    \item $p_x / p_y = 2.50$
    \item $p_x / p_y = 1.50$
    \item $p_x / p_y = 2.00$
    \item $p_x / p_y = 0.40$
\end{opts}

\q{Sean $U_A = 6\ln(x_A) + y_A$ y $U_B = 10\ln(x_B) + y_B$ con dotaciones $W_A = (2, 10)$ y $W_B = (6, 10)$. Fijando $p_y = 1$, el precio de equilibrio $p_x^*$ y las cantidades $(x_A^*, x_B^*)$ son:}
\begin{opts}
    \item $p_x^* = 1$ y $(x_A^*, x_B^*) = (4, 4)$
    \item $p_x^* = 2$ y $(x_A^*, x_B^*) = (3, 5)$
    \item $p_x^* = 2$ y $(x_A^*, x_B^*) = (4, 4)$
    \item $p_x^* = 4$ y $(x_A^*, x_B^*) = (2, 6)$
\end{opts}

\q{En una economía de intercambio puro con $U_A = x_A y_A$, $U_B = x_B y_B$ y dotación inicial $W_A = (9, 1)$ y $W_B = (1, 9)$, ¿qué intervalo de $x_A$ sobre la curva de contrato ($y_A = x_A$) define el Núcleo de la economía?}
\begin{opts}
    \item $x_A \in [1, 9]$
    \item $x_A \in [2, 8]$
    \item $x_A \in [3, 7]$
    \item $x_A \in [4, 6]$
\end{opts}

\q{En una economía con $U_A = x_A y_A$, $U_B = x_B y_B$ y dotaciones $W_A = (4, 2)$ y $W_B = (6, 8)$, se desea sostener la asignación eficiente $A = (7, 7)$ y $B = (3, 3)$ con $p_y = 1$. El precio $p_x^*$ y la transferencia de suma alzada $T$ de $B$ hacia $A$ son:}
\begin{opts}
    \item $p_x^* = 1$ y $T = 6$
    \item $p_x^* = 2$ y $T = 8$
    \item $p_x^* = 1$ y $T = 4$
    \item $p_x^* = 1$ y $T = 8$
\end{opts}

\q{Sean $U_A = \min(x_A, y_A)$ y $U_B = \min(x_B, y_B)$ con dotaciones $W_A = (5, 4)$ y $W_B = (5, 2)$. Normalizando $p_y = 1$, en el equilibrio walrasiano el precio $p_x^*$ y la utilidad de $A$ ($U_A^*$) son:}
\begin{opts}
    \item $p_x^* = 0$ y $U_A^* = 4$
    \item $p_x^* = 1$ y $U_A^* = 4$
    \item $p_x^* = 0$ y $U_A^* = 5$
    \item $p_x^* = 0.6$ y $U_A^* = 3$
\end{opts}

\q{Una economía cuenta con $\bar{L} = 90$ y $\bar{K} = 90$ para producir $X = L_X^{1/2} K_X^{1/2}$ e $Y = L_Y^{1/3} K_Y^{2/3}$. Si en una asignación técnicamente eficiente se emplea $L_X = 45$, ¿cuánto capital $K_X$ debe asignarse al sector $X$?}
\begin{opts}
    \item $K_X = 45$ unidades
    \item $K_X = 30$ unidades
    \item $K_X = 60$ unidades
    \item $K_X = 20$ unidades
\end{opts}

\q{Una economía tiene la FPP $X^2 + 4Y^2 = 200$ y un consumidor representativo con $U(X, Y) = X Y$. En el equilibrio general eficiente, la producción $(X^*, Y^*)$ y el precio relativo $p_X / p_Y$ son:}
\begin{opts}
    \item $(X^*, Y^*) = (5, 10)$ y $p_X / p_Y = 2.0$
    \item $(X^*, Y^*) = (10, 5)$ y $p_X / p_Y = 2.0$
    \item $(X^*, Y^*) = (10, 5)$ y $p_X / p_Y = 0.5$
    \item $(X^*, Y^*) = (8, 6)$ y $p_X / p_Y = 0.75$
\end{opts}

\q{En una economía con un solo factor $\bar{L} = 30$, las tecnologías son $X = 2L_X$ e $Y = 4L_Y$, y el consumidor representativo tiene $U = X^{1/2} Y^{1/2}$. En el óptimo global, $(X^*, Y^*)$ y $p_X / p_Y$ son:}
\begin{opts}
    \item $(X^*, Y^*) = (40, 40)$ y $p_X / p_Y = 1$
    \item $(X^*, Y^*) = (60, 30)$ y $p_X / p_Y = 0.5$
    \item $(X^*, Y^*) = (20, 80)$ y $p_X / p_Y = 4$
    \item $(X^*, Y^*) = (30, 60)$ y $p_X / p_Y = 2$
\end{opts}

\q{En una economía sobre su FPP con $CMg_X = 2X$, $CMg_Y = 4Y$ y $TMS_{X,Y} = Y/X$ común para todos los consumidores, el punto actual es $(X, Y) = (4, 8)$. Los valores de $TMT$ y $TMS$ y el ajuste eficiente son:}
\begin{opts}
    \item $TMT = 0.25$, $TMS = 2.0$; producir más $X$ y menos $Y$
    \item $TMT = 0.25$, $TMS = 2.0$; producir más $Y$ y menos $X$
    \item $TMT = 2.00$, $TMS = 0.25$; producir más $X$ y menos $Y$
    \item $TMT = 0.50$, $TMS = 2.00$; mantener la producción actual
\end{opts}

\q{En una economía de Robinson Crusoe con dotación de tiempo $L + l = 27$, utilidad $U(c, l) = c^{1/2} l^{1/2}$ y función de producción $c = 6L^{1/2}$, el trabajo óptimo $L^*$ y el consumo eficiente $c^*$ son:}
\begin{opts}
    \item $L^* = 16$ y $c^* = 24$
    \item $L^* = 9$ y $c^* = 18$
    \item $L^* = 9$ y $c^* = 27$
    \item $L^* = 18$ y $c^* = 18$
\end{opts}

\q{Una economía tiene FPP $X + 2Y = 120$ y consumidor representativo $U = X Y$. Si se aplica un impuesto ad-valorem de $100\%$ sobre el consumo de $X$ ($p_X^c = 2p_X^p$) devolviendo toda la recaudación como transferencia de suma alzada, el equilibrio distorsionado $(X_t, Y_t)$ es:}
\begin{opts}
    \item $(X_t, Y_t) = (60, 30)$
    \item $(X_t, Y_t) = (30, 45)$
    \item $(X_t, Y_t) = (40, 40)$
    \item $(X_t, Y_t) = (20, 50)$
\end{opts}

\q{En una economía $2\times 2$ de costos unitarios fijos con beneficio nulo, $p_X = 2w + r$ y $p_Y = w + 2r$. Si inicialmente $p_X = 10$ y $p_Y = 11$, y luego $p_X$ sube a $13$ manteniendo $p_Y = 11$, las nuevas remuneraciones $(w_1, r_1)$ son:}
\begin{opts}
    \item $w_1 = 4$ y $r_1 = 5$
    \item $w_1 = 6$ y $r_1 = 1$
    \item $w_1 = 5$ y $r_1 = 4$
    \item $w_1 = 5$ y $r_1 = 3$
\end{opts}

\q{En un modelo $2\times 2$ con pleno empleo y coeficientes fijos $2X + Y = \bar{L}$ y $X + 2Y = \bar{K}$, la dotación inicial $(\bar{L}, \bar{K}) = (40, 50)$ cambia a $(46, 50)$ a precios constantes. Las variaciones $(\Delta X, \Delta Y)$ son:}
\begin{opts}
    \item $\Delta X = +4$ y $\Delta Y = -2$
    \item $\Delta X = +3$ y $\Delta Y = +0$
    \item $\Delta X = +2$ y $\Delta Y = +2$
    \item $\Delta X = +6$ y $\Delta Y = -3$
\end{opts}

\q{Dada la Frontera de Posibilidades de Utilidad $U_A^2 + U_B^2 = 160$ y la función de bienestar social utilitarista ponderada $W = 3U_A + U_B$, la combinación óptima $(U_A^*, U_B^*)$ es:}
\begin{opts}
    \item $(U_A^*, U_B^*) = (10, 6)$
    \item $(U_A^*, U_B^*) = (12, 4)$
    \item $(U_A^*, U_B^*) = (8, 8)$
    \item $(U_A^*, U_B^*) = (4, 12)$
\end{opts}

\q{Considere una Frontera de Posibilidades de Utilidad lineal $2U_A + U_B = 120$ y una función de bienestar social de Nash $W = U_A \cdot U_B$. El óptimo social $(U_A^*, U_B^*)$ corresponde a:}
\begin{opts}
    \item $(U_A^*, U_B^*) = (40, 40)$
    \item $(U_A^*, U_B^*) = (60, 0)$
    \item $(U_A^*, U_B^*) = (30, 60)$
    \item $(U_A^*, U_B^*) = (20, 80)$
\end{opts}

\q{Si la Frontera de Posibilidades de Utilidad está dada por $U_A + 3U_B = 120$ y el planificador maximiza la función de bienestar social rawlsiana $W = \min(U_A, U_B)$, el óptimo social es:}
\begin{opts}
    \item $(U_A^*, U_B^*) = (60, 20)$
    \item $(U_A^*, U_B^*) = (40, 40)$
    \item $(U_A^*, U_B^*) = (90, 10)$
    \item $(U_A^*, U_B^*) = (30, 30)$
\end{opts}

\q{Dada la Frontera de Posibilidades de Utilidad $U_A + 4U_B = 100$ y la función de bienestar social CES $W = (U_A^{1/2} + U_B^{1/2})^2$, la asignación de utilidades que maximiza $W$ es:}
\begin{opts}
    \item $(U_A^*, U_B^*) = (80, 5)$
    \item $(U_A^*, U_B^*) = (20, 20)$
    \item $(U_A^*, U_B^*) = (60, 10)$
    \item $(U_A^*, U_B^*) = (40, 15)$
\end{opts}

\q{Un proyecto de infraestructura hídrica genera un beneficio neto de $+180$ millones para el grupo $A$ y un costo de $-110$ millones para el grupo $B$. ¿Qué rango de transferencia $T$ de $A$ hacia $B$ convierte el proyecto en una mejora estricta de Pareto?}
\begin{opts}
    \item $70 < T < 180$ millones
    \item $110 < T < 180$ millones
    \item $0 < T < 110$ millones
    \item $110 < T < 290$ millones
\end{opts}

\q{Un caudal disponible de $\bar{Q} = 60$ unidades de agua debe asignarse entre dos usuarios con beneficios marginales $BMg_1(q_1) = 100 - 2q_1$ y $BMg_2(q_2) = 85 - q_2$. La asignación eficiente $(q_1^*, q_2^*)$ y el precio sombra $\lambda^*$ son:}
\begin{opts}
    \item $(q_1^*, q_2^*) = (30, 30)$ y $\lambda^* = 40$
    \item $(q_1^*, q_2^*) = (35, 25)$ y $\lambda^* = 50$
    \item $(q_1^*, q_2^*) = (25, 35)$ y $\lambda^* = 50$
    \item $(q_1^*, q_2^*) = (20, 40)$ y $\lambda^* = 60$
\end{opts}

\q{Dos usuarios de una cuenca tienen $BMg_1 = 100 - 2q_1$ y $BMg_2 = 85 - q_2$ con un caudal total de $60$ unidades. Si parten de la asignación inicial $(q_1, q_2) = (45, 15)$, ¿cuál es la ganancia neta de bienestar al reasignar hasta el óptimo $(25, 35)$?}
\begin{opts}
    \item $300$ unidades monetarias
    \item $450$ unidades monetarias
    \item $900$ unidades monetarias
    \item $600$ unidades monetarias
\end{opts}

\q{Cinco hogares con preferencias unimodales votan por mayoría simple el aporte mensual para mejorar la red de agua local. Sus montos óptimos individuales son $S/\ 20$, $S/\ 35$, $S/\ 50$, $S/\ 80$ y $S/\ 120$. ¿Qué monto resulta ganador de Condorcet?}
\begin{opts}
    \item $S/\ 50$ mensuales
    \item $S/\ 61$ mensuales
    \item $S/\ 35$ mensuales
    \item $S/\ 80$ mensuales
\end{opts}

\q{En una economía Arrow-Debreu con dos estados equiprobables (húmedo $H$ y seco $S$), ambos agentes tienen $U_i = 0.5\ln(c_{iH}) + 0.5\ln(c_{iS})$ y dotaciones $W_A = (12, 4)$ y $W_B = (4, 4)$. El precio relativo del agua en estado seco frente al húmedo ($p_S / p_H$) es:}
\begin{opts}
    \item $p_S / p_H = 0.5$
    \item $p_S / p_H = 2.0$
    \item $p_S / p_H = 1.0$
    \item $p_S / p_H = 3.0$
\end{opts}

\q{Un hogar vulnerable con $U(x, y) = x^{1/2} y^{1/2}$ e ingreso $I = 72$ enfrenta $p_x = 9$ (agua) y $p_y = 1$. Si un subsidio tarifario reduce $p_x$ a $4$, ¿cuál es el costo fiscal del subsidio tarifario y qué transferencia de suma alzada daría al hogar exactamente la misma utilidad a precios sin distorsión?}
\begin{opts}
    \item Costo tarifario $= 45$ y suma alzada $= 45$
    \item Costo tarifario $= 36$ y suma alzada $= 27$
    \item Costo tarifario $= 45$ y suma alzada $= 36$
    \item Costo tarifario $= 54$ y suma alzada $= 36$
\end{opts}
